In this work we have developed the $r$-matrix formulation of integrable systems starting from the Lax pair description and progressively uncovering the underlying algebraic structure.

We first recalled that the existence of a Lax representation allows one to construct conserved quantities as spectral invariants of the Lax matrix. We then introduced the Poisson algebra of Lax matrices and showed that, when the Poisson brackets take the $r$-matrix form
\begin{equation}
    \{L_1,L_2\} = [r_{12},L_1] - [r_{21},L_2],
\end{equation}
the spectral invariants $\Tr(L^n)$ are in involution. We also discussed the converse statement, namely the local reconstruction of an $r$-matrix structure from the involution of the eigenvalues, under suitable regularity assumptions.

The algebraic content of this formalism was then developed through the tensor Casimir element and the classical and modified Yang--Baxter equations. In particular, we showed how the Yang--Baxter equation appears as the compatibility condition ensuring the Jacobi identity for the Poisson bracket, and how the Casimir element naturally enters the structure of the theory.

These general constructions were finally illustrated on the example of the open Toda chain associated with a simple Lie algebra. We reviewed its Lax representation and derived explicitly the corresponding $r$-matrix, showing that it reproduces the Poisson structure and yields commuting Hamiltonians. This provides a concrete realization of the abstract framework and clarifies how the general algebraic formalism is implemented in an integrable model.

For completeness, we also briefly reformulated the tensorial picture in terms of the corresponding $R$-matrix operator in the appendix.

Overall, the $r$-matrix formalism provides a unified perspective on integrability, linking Hamiltonian dynamics, spectral invariants, and Lie algebraic structures in a natural and explicit way.